\section{讨论与后续验证}
\label{sec:discussion}

本地结果提出一个明确的问题：在本次校准与搜索域内，所有候选均满足首词元 KL 上限，却没有候选达到关键词要求。它说明现有代理组合与预算没有产出可验收候选，不能定位原因究竟是冻结激活误差、末词元覆盖不足、目标尺度、搜索范围还是模型本身。下面将已有实现与仍需实验的假设分开。

\subsection{trajectory-v2：已实现但尚未实测}

当前审计版本 \code{2871bc1} 中，\code{ara\_trajectory.py} 沿基座生成的 continuation 做 teacher-forced 捕获，使用同生成步的正常与目标参考构建损失\cite{ara2026localv2}。模板每类使用 96 条提示、最多八个 continuation 位置，衰减系数 0.85，权重在每条提示内部归一化。相较 point-v1 单位置缓存，这扩展了被约束的位置，但仍然是冻结基座轨迹代理，编辑模型自由生成时可能偏离该轨迹。

v2 的 Gram 项是 $10^{-4}G(A,B)$，没有式~\eqref{eq:local-total} 的初始 Gram 分母，不能把两版目标写为相同归一化损失。\code{ara\_search.py} 与 \code{study\_runner.py} 实现八个搜索坐标，包括注意力与 MLP 部署增益，采用单 worker 的 8 个固定锚点、24 个 startup、88 个自适应试次调度。它不是 point-v1 的六维、36+84 协议，尚未测量扩大范围或增益带来的效果。

\code{scorers/refusal\_log\_odds.py} 提供连续的拒答与回答前缀对数几率分数；模板用该分数与 KL 作为优化目标，Keywords 保留为验收观察。该分数依赖预设前缀，不自动成为语义拒答真值。\code{protocol\_data.py} 约束校准、验证与审计角色；模板验证改用 \code{train[400:500]}，不能回填进 v1 的 \code{train[300:400]}。

验收相关模块 \path{acceptance.py}、\path{acceptance_export.py} 和 \path{artifact_schema.py} 实现候选选择、两次验证重放、导出前再次应用、单次审计账本、独立进程审计、指纹绑定与暂存后发布。运行保护还限定目标模块、设备、资源和部署漂移。这些是 E1 源码及配置事实；本次 v1 归档没有 v2 的候选、审计或导出结果，无法判断其是否提高成功率、语义效果或稳定性。

\subsection{从失败观察到可证伪假设}

第一项假设是冻结 I/O 的激活漂移限制了局部目标对整网行为的预测。可在相同预算下比较一次捕获、分块刷新与逐模块刷新，并直接测量缓存代理输出与编辑后真实输出之差。即使刷新降低代理误差，也仍需检查语义行为与能力指标，不能仅以局部损失证明改进。

第二项假设是末位置目标遗漏了较晚形成的拒答轨迹。point-v1 与 trajectory-v2 应使用匹配模型、提示与计算预算，比较首位置、多位置冻结轨迹及自由生成中的逐步行为。多位置成功与否需独立审计，不能从新增八个位置的代码推断。

第三项假设是当前搜索边界限制有效编辑。ID 66 的部分参数靠近上界是描述性现象，不证明最优解在域外。可预先指定范围扩展与部署增益消融，并与同域增加预算作配对比较。新一轮协议不得根据旧验证集反复选择后再宣称独立确认。

\subsection{统一比较、秩谱与资源}

单方向、SOM 多方向、上游全矩阵和本地低秩方法应固定模型修订、聊天模板、提示划分、目标模块、精度、解码与判别器\cite{arditi2024refusal,piras2026som}。预算同时报告候选数、生成次数、闭包调用、GPU 时长及峰值内存；独立种子需要分别运行，不能用同一搜索中的候选代替。验证集完成选择后，审计集只作预先约定的最终评价。HarmBench 可提供语义评估骨架\cite{mazeika2024harmbench}，仍需加入无害任务、长文本分布漂移和危险请求服从性。

组件消融应区分注意力、MLP 与二者联合编辑，并分别改变保持项、间隔惩罚和温度。硬 $k$NN 的 $k$ 扫描只适用于上游原型；本地软邻域不含该参数。更新奇异值谱、稳定秩及与既有方向的夹角可检验低秩更新的实际结构，但不能将几何相关性直接解释为拒答机制的数量或因果独立性。

完整资源报告需分别测量捕获、CPU 缓存、模块优化、验证与复载，保留失败和超时。上游的仅清零 $B$ 与完整重置、目标归一化与部署映射差异，可以作为历史实现消融；point-v1 已采用完整重置和无行归一化，不应再称这些为尚未实现的修复。

\subsection{行为解释与伦理边界}

拒答抑制既可帮助研究过度拒答与表征机制，也可能削弱原有防护。关键词减少不等于安全提升，危险知识是否仍可访问也不由局部几何决定。研究报告应同时呈现正常任务帮助性、危险请求服从性与失败统计。本文只提供公式、归档聚合和限制分析；本次验收失败没有产生可接受适配器，也没有证据将本地观察解释成更可靠的安全控制。
